\section{Edge/Fog Computing, Attribute Caching, dan Freshness}

\textit{Edge/fog computing} menjadi pendekatan penting dalam arsitektur IoT karena memindahkan sebagian proses komputasi dan pengambilan keputusan lebih dekat ke \textit{endpoint}. Dalam konteks kontrol akses, \textcite{kalaria_adaptive_2024} menunjukkan bahwa kontrol akses berbasis \textit{fog} dapat membawa \textit{decision-making} dan informasi kebijakan lebih dekat ke \textit{end nodes} sehingga kebijakan akses dapat dievaluasi secara lebih adaptif dan \textit{real-time}. Pendekatan ini juga mengurangi ketergantungan penuh pada \textit{cloud}, terutama ketika aplikasi IoT membutuhkan respons cepat dan konektivitas jaringan tidak selalu stabil. \textcite{wang_edge-enabled_2025} juga menekankan pentingnya autentikasi dan otorisasi berbasis \textit{edge} untuk mendukung latensi rendah dan manajemen akses luring pada layanan IoT berbasis \textit{edge}. Pada arsitektur seperti ini, \textit{gateway} atau \textit{fog node} dapat menjadi lokasi strategis untuk menempatkan komponen evaluasi dan penegakan kebijakan, seperti \textit{Policy Decision Point} (PDP) atau \textit{Policy Enforcement Point} (PEP) agar proses otorisasi tidak selalu dilakukan secara terpusat di \textit{cloud}.

Untuk meningkatkan responsivitas, atribut atau konteks yang sering digunakan dapat disimpan sementara pada \textit{edge} melalui mekanisme \textit{caching}. \textcite{weerasinghe_adaptive_2023} menunjukkan bahwa \textit{adaptive context caching} pada aplikasi IoT dapat meningkatkan efisiensi performa dan biaya ketika sistem harus merespons \textit{context queries} secara \textit{real-time}. Dalam konteks implementasi ABAC/XACML, \textit{caching} atribut dapat mengurangi kebutuhan pengambilan atribut berulang dari sumber informasi kebijakan, memperpendek waktu \textit{retrieval} atribut, dan menurunkan latensi evaluasi kebijakan. Namun, manfaat performa ini bergantung pada kemampuan sistem menjaga agar atribut yang tersimpan di \textit{cache} tetap cukup mutakhir untuk digunakan dalam keputusan akses.

\textit{Freshness} menjadi isu utama karena atribut yang tersimpan di \textit{cache} hanya valid dalam periode tertentu. \textcite{medvedev_refresh_2023} menjelaskan bahwa \textit{freshness, accuracy, completeness,} dan latensi merupakan bagian dari \textit{Quality of Service} (QoS) dalam pengelolaan konteks IoT sehingga strategi \textit{caching} perlu memperhitungkan \textit{trade-off} antara kualitas informasi dan biaya pengambilan data. Dalam praktik \textit{caching}, periode \textit{refresh} atau TTL yang lebih panjang dapat mengurangi frekuensi \textit{retrieval} atribut dan \textit{overhead} komunikasi sehingga meningkatkan efisiensi sistem. Sebaliknya, TTL yang lebih pendek dapat menjaga \textit{freshness} atribut karena data lebih sering diperbarui, tetapi meningkatkan frekuensi \textit{retrieval}, sinkronisasi, dan \textit{overhead} komunikasi yang dapat memengaruhi latensi evaluasi akses.

Masalah \textit{stale attribute} tidak hanya berdampak pada performa, tetapi juga pada kualitas dan keamanan keputusan otorisasi. \textcite{wang_edge-enabled_2025} menegaskan bahwa keputusan akses pada \textit{edge-IAM} menjadi dipertanyakan apabila sistem mengasumsikan nilai atribut selalu mutakhir dan berasal dari sumber yang sepenuhnya terpercaya. Penelitian tersebut juga mengusulkan LAA yang memasukkan \textit{freshness} nilai atribut dan \textit{trustworthiness} penyedia atribut agar kontrol akses menjadi lebih reliabel. Dengan demikian, atribut yang sudah usang dapat menyebabkan keputusan yang salah, baik berupa \textit{false allow} ketika akses yang seharusnya ditolak justru diizinkan, \textit{false deny} ketika akses sah ditolak, maupun \textit{revocation false allow} ketika hak akses yang sudah dicabut masih dianggap valid di \textit{edge}.

Oleh karena itu, \textit{caching} atribut pada \textit{edge} perlu diperlakukan sebagai persoalan optimasi keamanan, bukan sekadar optimasi performa. Strategi TTL yang seragam untuk semua atribut kurang memadai karena setiap atribut memiliki tingkat volatilitas dan dampak keamanan yang berbeda. Atribut yang sangat dinamis atau berdampak tinggi terhadap keamanan membutuhkan \textit{refresh} yang lebih ketat, sedangkan atribut yang lebih stabil dapat diberi TTL yang lebih panjang untuk mengurangi beban \textit{retrieval} atribut. Dengan demikian, konfigurasi TTL per atribut, ukuran \textit{cache}, dan pemilihan \textit{cacheable attribute} perlu dioptimalkan agar sistem dapat menjaga keseimbangkan antara otorisasi berlatensi rendah, efisiensi komunikasi, \textit{freshness} atribut, dan kualitas keputusan akses.